\originalchapter{分子的转动与振动}\label{ch:molecule}
\originalsection{自由度与正常模式}
双原子气体在室温的定容热容约为每分子$5\kb/2$, 更低温趋$3\kb/2$, 高温振动激活后趋$7\kb/2$. 分子没有突然失去自由度; 改变的是热能是否足以占据相应激发态.

含$s$个原子的稳定分子共有$3s$个位移坐标. 质心平移占3个, 线性分子独立转动占2个, 非线性分子占3个, 剩余分别为$3s-5$或$3s-6$个振动正常模式. 单原子没有内部转动或振动. 在平衡构型附近展开势能$V=V_0+\frac12\sum_{ij}K_{ij}u_iu_j+O(u^3)$, 用质量加权坐标将动能化为单位质量二次型, 再正交对角化刚度矩阵. 零特征值来自刚体平移和转动的对称性, 非零正特征值给振动频率平方. 稳定允许零模, 所以“正半定”比“正定”准确.

忽略分子间作用与转动振动耦合时$z=z_{\rm cm}z_{\rm rot}z_{\rm vib}$, 各部分自由能和热容相加. 平移通常可作经典近似, 但转动、振动能隙可能大于$\kb T$. 完全经典的热容是每分子$\kb[3/2+r/2+(3s-3-r)]$, $r$为独立转动数. 例如线性CO$_2$有4振动模式, 非线性H$_2$O有3振动模式; 退化模式也要按维数逐个计数.

\originalsection{量子谐振子}
单模能量$E_n=\hbar\omega(n+1/2)$, $n\ge0$. 几何求和给
\[
z_{\rm vib}=\frac{\e^{-x/2}}{1-\e^{-x}},\quad x=\beta\hbar\omega,\qquad
\avg E=\frac{\hbar\omega}2+\frac{\hbar\omega}{\e^x-1}.
\]
振动量子数概率是$(1-\e^{-x})\e^{-nx}$, 与几何分布相同. 热容
\[
c_{vib}=\kb\frac{x^2\e^x}{(\e^x-1)^2}
\]
在$x\ll1$趋$\kb$, 在$x\gg1$为$\kb x^2\e^{-x}$. 低温主要是基态和第一激发态之间的罕见转换, 因而能隙控制指数; 高温可以把稠密能级和近似为积分, 恢复经典结果. 零点能不贡献固定频率的热容, 但若频率随体积变, 它会贡献压强, 不能在所有问题中任意丢掉.

\originalsection{量子转子与光谱}
线性分子绕垂直轴的转动惯量为$I$, 能量$\epsilon_l=\hbar^2l(l+1)/(2I)=\kb\Theta_Rl(l+1)$, 每能级的磁量子数$m=-l,\ldots,l$使简并度$2l+1$. 配分函数
\[
z_R=\sum_{l=0}^{\infty}(2l+1)\e^{-\Theta_Rl(l+1)/T}.
\]
低温只留$l=0,1$, $z_R\simeq1+3\e^{-2\Theta_R/T}$, 于是$c_R\simeq3\kb(2\Theta_R/T)^2\e^{-2\Theta_R/T}$. 高温用变量$u=l(l+1)\Theta_R/T$将积分化为$z_R\simeq T/\Theta_R$, 给平均转动能$\kb T$和热容$\kb$. 这对应两个二次动量自由度. 同核分子的交换对称性与核自旋会限制允许$l$, 不能把这个未限制转子式直接当作所有H$_2$的精确模型.

Raman散射满足分子能量增量$\Delta\epsilon=h(\nu_i-\nu_f)$. 定义光频位移 $\Delta\nu=\nu_f-\nu_i$，转动选择规则$\Delta l=\pm2$使Stokes位移为$-(4l+6)\kb\Theta_R/h$, 相邻线间距$4\kb\Theta_R/h$. 强度在忽略态依赖散射矩阵元的近似下与初态分子数$(2l+1)\e^{-\Theta_Rl(l+1)/T}$成正比. 它反映简并度与热权重两方面, 并不是单纯读取能级间隔.

两个相同核交换使转动空间波函数获得$(-1)^l$. 对H$_2$的质子是自旋$1/2$Fermion: 核自旋单重态反对称, 配偶数$l$的对称空间态, 称para; 三重态对称, 配奇数$l$, 称ortho. 相应转动配分函数是$z_p=\sum_{l\ {even}}(2l+1)\e^{-u l(l+1)}$和$z_o=3\sum_{l\ {odd}}(2l+1)\e^{-u l(l+1)}$, $u=\Theta_R/T$. 平衡混合使用$z_p+z_o$, 高温ortho:para趋3:1. 现实转换慢时, 比例可被历史制备冻结, 就需固定两种分子数分别求热容. D$_2$的自旋1核为Boson, 核自旋对称态有6个、反对称态有3个, 因而偶奇$l$强度比反过来是6:3. 这些统计权重解释了交替谱线强度, 不必假设分子间有新相互作用.
